\subsection{Исполнительная подсистема}\label{subsec:ch1/sec2/sub3}

Любая машина, в том числе и НТТС, состоит из отдельных частей: силового (приводы) и рабочего оборудования, передаточных устройств, системы управления и других.

Исполнительная подсистема (ИcП) является ключевым элементом автоматизации НТТС, обеспечивая непосредственное выполнение управляющих команд, сформированных системой управления. Основная задача ИcП заключается в преобразовании электрических сигналов в механические действия, необходимые для перемещения и выполнения рабочих операций НТТС. Объектами роботизации (автоматизации) в строительстве являются технологические операции, выполняемые при возведении зданий и сооружений. Например процессы земляные, монтажные или отделочные.

\nomenclature{ИcП}{исполнительная подсистема\nomrefpage}

В состав ИcП входят следующие основные компоненты \cite{Vilman, Evtukov}:

\begin{itemize}
	\item \textit{приводы}:
	устройства, преобразующие электрическую энергию в механическое движение. В зависимости от типа НТТС и решаемых задач могут использоваться электродвигатели, гидравлические или пневматические приводы. Электродвигатели обеспечивают высокую точность и быстродействие, в то время как гидравлические и пневматические системы применяются для задач, требующих значительных усилий, например, при управлении тяжелыми механизмами или подъемными устройствами;

	\item \textit{рабочие органы}:
	представляют собой рабочую часть манипулятора, служат для непосредственного выполнения технических операций или вспомогательных переходов. Включают в себя механизмы, отвечающие за выполнение специфических операций, таких как подъем, перемещение или обработка грузов. Например, в случае с автономными погрузчиками или экскаваторами это могут быть гидравлические манипуляторы, ковши или другие рабочие инструменты ;

	\item \textit{тормозные системы}:
	обеспечивают управление скоростью и остановку НТТС. В автоматизированных системах тормозные механизмы могут быть интегрированы с датчиками обратной связи, что позволяет точно контролировать процесс торможения в зависимости от дорожных условий и текущей скорости;

	\item \textit{рулевые механизмы}:
	отвечают за изменение направления движения НТТС. В автоматизированных системах рулевое управление реализуется с помощью электромеханических или гидравлических усилителей, которые обеспечивают точное выполнение команд по изменению траектории движения;

	\item \textit{системы обратной связи}:
	обеспечивают контроль за состоянием исполнительных устройств и передачу данных о выполнении команд в систему управления. Это позволяет корректировать управляющие воздействия в~реальном времени, минимизируя ошибки и повышая точность выполнения операций.
\end{itemize}

Особенностью ИcП НТТС является необходимость обеспечения высокой надежности и отказоустойчивости, что особенно важно в условиях эксплуатации, связанных с повышенными нагрузками и вибрациями. Для этого применяются резервирование критически важных компонентов, использование защищенных от внешних воздействий конструкций, а также интеграция систем диагностики, способных своевременно выявлять и устранять неисправности.

Кроме того, ИcП должна быть адаптирована к специфике работы НТТС, включая возможность управления сложными конструкциями, такими как полуприцепы или сменное рабочее оборудование (рис.\:\cref{fig:manipul}). Это требует разработки специализированных алгоритмов управления, учитывающих изменение динамических характеристик и массово-инерционных параметров системы \cite{PevznerUgaiSulemanov2010}.

\begin{figure}[ht]
	\centering{
		\includegraphics[scale=0.5]{views/manipulator.png}
	}
	\caption{Пример полуприцепа и рабочего оборудования}\label{fig:manipul}
\end{figure}

Одновременно с этим, являясь ключевым элементом автоматизации НТТС, ИcП может оказывать значительное влияние на другие подсистемы, как положительное, так и отрицательное. Это влияние может проявляться в виде ограничений, накладываемых на работу других подсистем, или в виде необходимости адаптации их функциональности для обеспечения согласованной работы всей системы в целом. Рассмотрим основные аспекты этого влияния.

ИcП может создавать дополнительные шумы и вибрации, которые могут негативно сказаться на точности работы датчиков, особенно чувствительных к механическим воздействиям. Например, вибрации от работы двигателей или гидравлических систем могут искажать данные акселерометров и гироскопов, что затрудняет точное определение положения и ориентации НТТС.

Электромагнитные помехи от мощных электродвигателей или силовых кабелей могут влиять на работу радаров, лидаров и других электронных датчиков \cite{confbib1}.

Для минимизации этих эффектов требуется тщательное проектирование системы, включая экранирование датчиков, использование фильтров для подавления шумов и вибраций, а также алгоритмов компенсации помех.

ИcП может влиять на точность позиционирования и локализации, особенно в динамических условиях: изменение массы и центра тяжести при работе с грузами или сменным оборудованием (например, полуприцепами) может привести к смещению расчетных параметров движения, что требует корректировки алгоритмов локализации.

Задержки в выполнении команд (например, из-за инерции тяжелых механизмов) могут привести к рассинхронизации между расчетной и фактической траекторией движения, что особенно критично при работе в условиях ограниченного пространства.

Для устранения этих проблем необходимо учитывать динамические характеристики ИcП при разработке алгоритмов управления и локализации, а также использовать системы обратной связи для оперативной корректировки.

ИcП может накладывать ограничения на работу подсистемы управления, особенно в части быстродействия -- механические системы (например, гидравлические приводы) могут иметь ограниченную скорость реакции, что требует адаптации алгоритмов управления для учета этих задержек. Энергопотребления -- мощные исполнительные устройства (например, двигатели или гидравлические насосы) могут потреблять значительное количество энергии, что требует оптимизации работы других подсистем для обеспечения баланса энергопотребления. Надежности -- отказ исполнительной подсистемы (например, заклинивание рулевого механизма или потеря управления тормозами) может привести к критическим последствиям, что требует разработки резервных контуров управления и систем диагностики.

ИcП может создавать дополнительные ограничения при работе в сложных условиях:

\begin{itemize}
	\item \textit{по маневренности}: крупные габариты и сложная динамика НТТС (особенно при наличии полуприцепов или сменного оборудования) могут затруднять движение в узких пространствах или на пересеченной местности;

	\item \textit{по скорости}: необходимость обеспечения безопасности и устойчивости НТТС может требовать снижения скорости движения, что влияет на общую производительность системы.
\end{itemize}

Для снижения негативного влияния ИcП на другие подсистемы могут быть применены следующие меры:

\begin{itemize}
	\item \textit{интеграция систем обратной связи}: для оперативного учета изменений в работе исполнительных устройств и корректировки управляющих команд;

	\item \textit{оптимизация алгоритмов управления}: с учетом динамических характеристик ИcП и возможных задержек;

	\item \textit{резервирование}: для повышения надежности и отказоустойчивости системы;

	\item \textit{адаптация под условия эксплуатации}: например, использование специализированных алгоритмов для работы с полуприцепами или в условиях пересеченной местности.
\end{itemize}

Таким образом, ИcП, хотя и является важным элементом автоматизации НТТС, может накладывать значительные ограничения на работу других подсистем. Однако при грамотном проектировании и интеграции этих ограничений можно минимизировать, обеспечивая согласованную и эффективную работу всей системы в целом.